\section*{Глава 3. Разработка веб-приложения}
\addcontentsline{toc}{section}{Глава 3. Разработка веб-приложения}
\stepcounter{section}

\subsection{Разработка базовой структуры приложения и вёрстка шаблонов страниц}

Проект организован в соответствии с типовой структурой Flask-приложения. Ключевые файлы:

\begin{itemize}
    \item \texttt{app.py} — создание и настройка приложения, все маршруты;
    \item \texttt{models.py} — модели базы данных (SQLAlchemy);
    \item \texttt{forms.py} — WTForms для валидации пользовательского ввода;
    \item папка \texttt{templates/} — HTML-шаблоны (Jinja2);
    \item папка \texttt{static/} — статические файлы (CSS, JavaScript).
\end{itemize}

Базовый шаблон \texttt{base.html} содержит навигационную панель, блок для flash-сообщений и 
подключение стилей Bootstrap (тема Minty). Навигационная панель динамически меняется: для 
неавторизованных пользователей показываются ссылки «Войти» и «Регистрация», для авторизованных
 — «Добавить цитату», «Импорт», «Экспорт», «Аналитика» и «Выйти». Если пользователь имеет 
роль \texttt{admin}, дополнительно отображается ссылка «Админ-панель».

Все остальные шаблоны наследуют \texttt{base.html} и переопределяют блок \texttt{content}. 
Для стилизации используются компоненты Bootstrap: карточки для цитат, формы для ввода данных, 
модальные окна для подтверждения действий. Адаптивная вёрстка обеспечивает корректное 
отображение на экранах разной ширины.

\subsection{Реализация аутентификации и авторизации с разграничением прав доступа}

Для управления аутентификацией используется расширение Flask-Login. В файле \texttt{models.py} 
определена модель \texttt{User}, которая наследует \texttt{UserMixin} — это даёт готовые методы 
для работы с сессиями. Модель содержит поля: \texttt{id} (первичный ключ), \texttt{login} 
(уникальное имя пользователя), \texttt{password\_hash} (хеш пароля), \texttt{role} 
(роль: \texttt{user} или \texttt{admin}), \texttt{created\_at} (дата регистрации). 
Листинг~\ref{lst:user_model} показывает модель пользователя.

\begin{lstlisting}[caption={Модель пользователя (models.py)}, label={lst:user_model}]
class User(UserMixin, db.Model):
    __tablename__ = 'users'
    id = db.Column(db.Integer, primary_key=True)
    login = db.Column(db.String(100), unique=True, nullable=False)
    password_hash = db.Column(db.String(256), nullable=False)
    role = db.Column(db.String(20), nullable=False, default='user')
    created_at = db.Column(db.DateTime, default=datetime.now)
\end{lstlisting}

Хеширование паролей выполняется при помощи функций \texttt{generate\_\-password\_hash} и 
\texttt{check\_password\_hash} из библиотеки \texttt{werkzeug.security}. При регистрации 
пароль хешируется, после чего хеш сохраняется в базу данных. При входе введённый пароль 
хешируется повторно и сравнивается с сохранённым хешем.

Маршрут \texttt{/register} обрабатывает регистрацию нового пользователя. После успешного 
создания записи в базе данных выполняется автоматический вход с помощью 
\texttt{login\_user(user)} — это избавляет пользователя от необходимости вводить логин и 
пароль сразу после регистрации. Маршрут \texttt{/login} проверяет логин и пароль, при 
успешной проверке вызывает \texttt{login\_user}. Маршрут \texttt{/logout} завершает сессию 
через \texttt{logout\_user}. Для защиты маршрутов, требующих авторизации, используется 
декоратор \texttt{@login\_required}.

Для разграничения прав доступа создан декоратор \texttt{@admin\_required}. Он проверяет, 
что пользователь авторизован и имеет роль \texttt{admin}. Если проверка не пройдена, 
пользователь получает сообщение об ошибке и перенаправляется на главную страницу.

\begin{lstlisting}[caption={Декоратор для проверки прав администратора (app.py)}, label={lst:admin_decorator}]
def admin_required(f):
    @wraps(f)
    def decorated_function(*args, **kwargs):
        if not current_user.is_authenticated or current_user.role != 'admin':
            return redirect(url_for('index'))
        return f(*args, **kwargs)
    return decorated_function
\end{lstlisting}

Все маршруты, относящиеся к администрированию (управление фильмами, управление жанрами), 
защищены декораторами \texttt{@login\_required} и \texttt{@admin\_required}.

\subsection{Реализация CRUD-интерфейса для сущностей: фильмы (общий каталог), цитаты, пользователи}

В приложении реализованы модели \texttt{User}, \texttt{Genre}, \texttt{Movie} и 
\texttt{Quote}. Связи между ними: фильм принадлежит одному жанру, цитата принадлежит 
одному фильму и одному пользователю. При удалении фильма все связанные с ним цитаты 
удаляются автоматически. Листинг~\ref{lst:quote_model} показывает модель \texttt{Quote}.

\begin{lstlisting}[caption={Модель цитаты (models.py)}, label={lst:quote_model}]
class Quote(db.Model):
    __tablename__ = 'quotes'
    id = db.Column(db.Integer, primary_key=True)
    text = db.Column(db.Text, nullable=False)
    character_name = db.Column(db.String(255))
    timestamp = db.Column(db.String(20))
    movie_id = db.Column(db.Integer, db.ForeignKey('movies.id', ondelete='CASCADE'), nullable=False)
    user_id = db.Column(db.Integer, db.ForeignKey('users.id', ondelete='CASCADE'), nullable=False)
    created_at = db.Column(db.DateTime, default=datetime.now)
\end{lstlisting}

\textbf{Создание цитаты.} Форма добавления цитаты содержит выпадающий список фильмов, 
поле для текста, поля для персонажа и таймкода. Выпадающий список заполняется из таблицы 
\texttt{movies}. При отправке формы создаётся объект \texttt{Quote}, поля заполняются 
данными из формы, а \texttt{user\_id} подставляется из \texttt{current\_user.id}. Объект 
добавляется в сессию SQLAlchemy и сохраняется в базу через \texttt{db.session.commit()}. 

\textbf{Чтение цитат.} Главная страница выполняет запрос 
\texttt{\texttt{Quote.query.filter\-\_by(...).order\_by(...).all()}}
который возвращает все цитаты текущего пользователя, отсортированные по дате добавления 
(сначала новые). Для детального просмотра используется маршрут \texttt{/quote/<int:id>}, 
который извлекает цитату по ID с помощью \texttt{get\_or\_404}.

\textbf{Обновление цитаты.} Форма редактирования предварительно заполняется текущими 
значениями цитаты (параметр \texttt{obj=quote} при создании формы). При отправке формы 
значения из формы присваиваются объекту \texttt{quote}, после чего изменения фиксируются 
в базе. Доступ к редактированию имеют только автор цитаты или администратор 
(проверка \texttt{quote.user\_id != current\_user.id and current\_user.role != 'admin'}).

\textbf{Удаление цитаты.} При переходе по маршруту \texttt{/quote/delete/<int:id>} объект 
цитаты удаляется из сессии с помощью \texttt{db.session.delete(quote)}. Перед удалением 
выполняется проверка прав. Удаление происходит после подтверждения через модальное окно.

\textbf{Фильмы (общий каталог).} Фильмы могут добавлять все пользователи (через AJAX-форму 
на странице создания цитаты). При добавлении проверяется, существует ли фильм с таким 
названием; если нет — создаётся новый. Администратор имеет доступ к странице 
\texttt{/admin/movies}, где отображаются все фильмы. Редактирование фильма выполняется 
через форму, удаление — через подтверждение. При удалении фильма все связанные с ним цитаты 
удаляются автоматически.

\textbf{Жанры.} Жанры управляются только администратором. При добавлении фильма пользователь 
выбирает жанр из выпадающего списка, который заполняется из таблицы \texttt{genres}. 
Администратор может добавлять новые жанры, редактировать названия и удалять жанры. При 
удалении жанра поле \texttt{genre\_id} у связанных фильмов становится пустым (\texttt{NULL}).

\textbf{AJAX-добавление фильма.} На странице создания цитаты встроена форма добавления 
фильма. При нажатии кнопки «Сохранить фильм» отправляется AJAX-запрос на сервер. Сервер 
создаёт фильм и возвращает JSON с его ID и названием. На клиенте новый фильм добавляется 
в выпадающий список и автоматически выбирается. Это позволяет пользователю не покидать 
страницу для добавления нового фильма. AJAX-маршрут для добавления фильма прелставлен в 
листинге 4.

\newpage
\begin{lstlisting}[caption={AJAX-добавление фильма (app.py)}, label={lst:ajax_movie}]
@app.route('/movie/add/ajax', methods=['POST'])
@login_required
def add_movie_ajax():
    data = request.get_json()
    title = data.get('title')
    if not title:
        return jsonify({'success': False, 'error': 'Название фильма обязательно'})
    genre_id = data.get('genre_id')
    if not genre_id:
        return jsonify({'success': False, 'error': 'Выберите жанр'})
    movie = Movie(
        title=title,
        director=data.get('director') or None,
        release_year=data.get('release_year') or None,
        genre_id=int(genre_id),
        added_by=current_user.id
    )
    db.session.add(movie)
    db.session.commit()
    return jsonify({'success': True, 'movie_id': movie.id, 'movie_title': movie.title})
\end{lstlisting}

\subsection{Реализация агрегирования данных и аналитики: личная статистика пользователя}

Страница аналитики показывает пользователю три графика, построенных на основе его цитат: 
топ-5 фильмов по количеству цитат, распределение цитат по жанрам и динамику добавления цитат 
по месяцам.

Данные для графиков формируются с помощью агрегирующих SQL-запросов. Для топ-5 фильмов 
выполняется запрос, который соединяет таблицы \texttt{movies} и \texttt{quotes}, фильтрует 
цитаты текущего пользователя, группирует записи по фильмам, подсчитывает количество цитат 
в каждой группе, сортирует по убыванию и ограничивает результат пятью записями:

\newpage
\begin{lstlisting}[caption={SQLAlchemy-запрос для топ-5 фильмов пользователя}, label={lst:top_movies}]
user_top_movies = db.session.query(
    Movie.title,
    func.count(Quote.id).label('count')
).join(Quote).filter(
    Quote.user_id == current_user.id
).group_by(Movie.id).order_by(
    func.count(Quote.id).desc()
).limit(5).all()
\end{lstlisting}

Для распределения по жанрам выполняется двойное соединение (через таблицу \texttt{movies} к 
таблице \texttt{genres}), группировка по названию жанра и подсчёт цитат. Для динамики по 
месяцам используется функция \texttt{extract('month', Quote.created\_at)}, которая извлекает 
номер месяца из даты; затем выполняется группировка по месяцу и подсчёт цитат. После 
выполнения запроса создаётся массив из 12 элементов (по числу месяцев), который заполняется 
полученными значениями — это гарантирует, что на графике будут отображены все месяцы, даже 
если в некоторые из них цитаты не добавлялись.

После выполнения запросов данные передаются в шаблон. Для визуализации используется 
библиотека Chart.js. На странице создаются три элемента \texttt{canvas}, для каждого 
инициализируется график соответствующего типа: \texttt{bar} для гистограммы (топ фильмов), 
\texttt{pie} для круговой диаграммы (жанры), \texttt{line} для линейного графика (динамика 
по месяцам). Для линейного графика параметр \texttt{tension} установлен в 0.1, что делает 
линию достаточно резкой для наглядного отображения изменений.

Для администратора на странице аналитики доступно переключение между личной и общей 
статистикой. Общая аналитика строится по тем же принципам, но без фильтрации по 
\texttt{user\_id} (учитываются цитаты всех пользователей). Дополнительно отображается 
график топ-5 активных пользователей, который формируется запросом с группировкой по 
\texttt{user\_id} и подсчётом цитат.

\subsection{Реализация импорта цитат из файлов (CSV/JSON) и экспорта в CSV/JSON}

Импорт цитат выполняется через загрузку CSV или JSON-файла. На странице \texttt{import.html} 
находится форма с полем выбора файла. После загрузки файла сервер читает его содержимое, 
декодирует в UTF-8 и определяет формат по расширению.

Для CSV-файлов используется модуль \texttt{csv} и класс \texttt{DictReader}, который 
преобразует каждую строку в словарь, где ключи — это заголовки столбцов. Ожидаемые 
заголовки: «фильм», «цитата», «персонаж» (опционально), «таймкод» (опционально). Для 
каждой строки проверяется, существует ли фильм с указанным названием; если нет — фильм 
создаётся автоматически. Затем проверяется, нет ли уже у пользователя цитаты с таким же 
текстом и таким же фильмом — если есть, строка пропускается (защита от дубликатов). После 
обработки всех строк пользователь получает сообщение с количеством добавленных и пропущенных 
цитат.

Для JSON-файлов используется модуль \texttt{json}. Ожидается массив объектов с полями 
«фильм», «цитата», «персонаж», «таймкод». Логика импорта аналогична CSV.

Экспорт цитат реализован через маршруты \texttt{/export/csv} и \texttt{/export/json}. 
Сервер формирует файл с цитатами текущего пользователя и отправляет его на скачивание. 
Для CSV используется объект \texttt{StringIO} (хранение данных в памяти без создания файла 
на диске) и \texttt{csv.writer}. Первая строка — заголовки. Для каждой цитаты записывается 
строка с названием фильма, текстом цитаты, персонажем, таймкодом и датой добавления. Для 
JSON формируется список словарей, который преобразуется в JSON-строку с помощью 
\texttt{json.dumps} (параметр \texttt{ensure\_ascii=False} сохраняет русские буквы). 
В обоих случаях в HTTP-ответе устанавливается заголовок \texttt{Content-Disposition: 
attachment}, чтобы браузер скачивал файл, а не отображал его.

\begin{lstlisting}[caption={Экспорт цитат в CSV (app.py)}, label={lst:export_csv}]
@app.route('/export/csv')
@login_required
def export_csv():
    quotes = Quote.query.filter_by(user_id=current_user.id).all()
    si = StringIO()
    writer = csv.writer(si)
    writer.writerow(['Фильм', 'Цитата', 'Персонаж', 'Таймкод', 'Дата добавления'])
    for quote in quotes:
        writer.writerow([
            quote.movie.title if quote.movie else '',
            quote.text,
            quote.character_name or '',
            quote.timestamp or '',
            quote.created_at.strftime('%d.%m.%Y %H:%M')
        ])
    output = si.getvalue().encode('utf-8-sig')
    return Response(output, mimetype='text/csv',
                    headers={'Content-Disposition': 'attachment; filename=quotes.csv'})
\end{lstlisting}